% =====================================================================
% 01_introduction.tex
% =====================================================================
\begingroup
\newcolumntype{L}[1]{>{\hsize=#1\hsize\raggedright\arraybackslash}X}

\section{Introduction}
\label{sec:intro}

\subsection{What Is a WLAN?}

A \ac{wlan} is a data network that connects devices over a radio-frequency
medium instead of cables, typically inside a home, an office floor, a campus,
or a public venue. Modern \acp{wlan} are almost exclusively built on the \ac{ieee}
802.11 family of standards, which define the \ac{phy} and \ac{mac} layers
\cite{ieee80211_2020}.

\subsection{LAN and WLAN}

A LAN connects devices within a limited area, usually using switched Ethernet
and physical cables. A WLAN provides the same local-network function over a
shared wireless medium, giving users mobility and easier deployment but making
performance more sensitive to distance, interference and obstacles. A WLAN
normally connects to the rest of the LAN through an access point.

\subsection{WLAN Compared with Other Wireless Networks}

Wireless networks are commonly classified by coverage area
(Table~\ref{tab:wireless-classes}). \acp{wlan} sit between short-range personal
networks and wide-area cellular networks.

\begin{table}[htbp]
  \centering\footnotesize
  \caption{Classes of wireless networks (typical values).}
  \label{tab:wireless-classes}
  \begin{tabularx}{\linewidth}{@{} L{0.8} L{0.9} L{0.9} L{1.2} L{1.2} @{}}
    \toprule
    \textbf{Class} & \textbf{Typical range} & \textbf{Examples} & \textbf{Spectrum}        & \textbf{Typical use}         \\
    \midrule
    \ac{wpan}      & up to about 10~m       & Bluetooth         & Unlicensed (2.4~GHz)     & Headsets, sensors, wearables \\
    \acs{wlan}     & tens of metres indoors & Wi-Fi (802.11)    & Unlicensed 2.4, 5, 6~GHz & Homes, offices, campuses     \\
    Wireless MAN   & several kilometres     & WiMAX (802.16)    & Licensed and unlicensed  & Fixed broadband, backhaul    \\
    \ac{wwan}      & kilometres per cell    & LTE               & Licensed                 & Mobile broadband, voice      \\
    \bottomrule
  \end{tabularx}
\end{table}

\subsection{Advantages and Limitations}

In normal use, a station associates with an access point and exchanges frames
over the shared wireless medium. This model is used in homes, offices, campuses,
public hotspots and temporary networks where mobility or avoiding new cabling is
important.

\begin{itemize}
  \item \textbf{Advantages:}
    \begin{itemize}
      \item \textbf{Mobility and convenience:} users keep connectivity while
        moving between rooms or buildings.
      \item \textbf{Fast, low-cost deployment:} no cable pulling; ideal for
        temporary sites, heritage buildings and dense venues.
      \item \textbf{Scalability:} capacity grows by adding \acp{ap} and
        channels.
    \end{itemize}
  \item \textbf{Disadvantages:}
    \begin{itemize}
      \item \textbf{Shared medium:} all stations in range of the same
        channel compete for the same airtime, so capacity is lower and
        less predictable than with switched Ethernet.
      \item \textbf{Interference and propagation effects:} walls,
        reflections and other radios degrade performance.
      \item \textbf{Security exposure:} radio signals leave the building, so
        strong encryption and authentication are mandatory.
    \end{itemize}
\end{itemize}

\begin{tomtat}
\item \acs{wlan} là mạng cục bộ dùng sóng vô tuyến thay cho cáp, hầu hết xây dựng trên họ chuẩn \acs{ieee} 802.11 (định nghĩa lớp \acs{phy} và \acs{mac}).
\item So với \acs{lan} có dây, \acs{wlan} linh hoạt và dễ triển khai hơn nhưng dùng chung môi trường truyền, nên nhạy với khoảng cách, nhiễu và vật cản.
\item Về phạm vi, \acs{wlan} nằm giữa \acs{wpan} (khoảng 10~m) và \acs{wwan} (hàng km mỗi ô).
\item Di động, dễ triển khai nhanh và rẻ, mở rộng bằng cách thêm \acs{ap} và kênh. Nhưng có ít dung lượng để chia sẻ, dễ nhiễu, sóng lan ra diện rộng nên bắt buộc phải có mã hóa và xác thực mạnh.
\end{tomtat}
\endgroup
